\chapter{Lean inventory: declarations and mathematical meaning}

\section{How to read the formalization}

The repository uses three kinds of declarations:
\begin{enumerate}
\item \textbf{definitions}: mathematical objects and predicates;
\item \textbf{theorems}: propositions proved from the current Lean
      dependencies and earlier declarations;
\item \textbf{structures}: explicit data packages whose existence is still
      an input unless a constructor theorem is provided.
\end{enumerate}
The distinction is essential.  A structure such as
\lean{RootInput} or \lean{PhaseParameterPuzzleData} is not an axiom merely
because a theorem accepts it as an argument.

\section{Public roots and exact Mandelbrot approximations}

\begin{description}
\item[\lean{Mlc/Core.lean}]
Defines \lean{MLC.mlc_conjecture}: the compatibility theorem from
\lean{RootInput} to \lean{LocallyConnectedSpace M}.  Its proof uses only the
categorical root and the foundation axioms reported by the audit.

\item[\lean{Mlc/CategoricalRoot.lean}]
Defines \lean{MLC.RootInput},
\lean{Categorical.mandelbrotObject}, and
\lean{Categorical.MLCConjecture}.  The theorem
\lean{Categorical.mlcConjecture_iff_set} is definitional equivalence between
the categorical and subspace targets.

\item[\lean{Mlc/CategoricalMandelbrot.lean}]
Defines the increasing inner and decreasing outer orbit systems.  The
theorems
\lean{Categorical.Mandelbrot.iUnion_innerOrbitSet_eq_set} and
\lean{Categorical.Mandelbrot.iInter_outerOrbitSet_eq_set} prove the two exact
limit identities.  Compactness is
\lean{Categorical.Mandelbrot.isCompact_outerOrbitSet}.
\end{description}

\section{Topology and local-connectivity bridges}

\begin{description}
\item[\lean{Mlc/LocalConnectivity.lean}]
Transfers connectedness between a subspace and its image and proves the
connected-basis lemma used by parameter puzzles.

\item[\lean{Mlc/ParameterComponentApproximation.lean}]
Defines \lean{ParameterComponent.UniformOuterBuffer} and proves
\lean{ParameterComponent.locallyConnectedSpace_of_uniformOuterBuffer}.
This is the set-level theorem behind the categorical root.

\item[\lean{Mlc/FiniteComponentCriterion.lean}]
Defines rational local pieces and finite cover certificates.  The theorem
\lean{FiniteComponent.uniformOuterBuffer_of_finiteComponentCriterion}
turns finite certificates into the uniform buffer.

\item[\lean{Mlc/RetractionLocalConnectivity.lean}]
Proves that a continuous retraction from a convex ambient domain onto
$\M$ yields local connectivity.

\item[\lean{Mlc/FlowInterfaces.lean}]
Defines \lean{FlowInterfaces.CoherentParametrization},
\lean{FlowInterfaces.TerminalExtension}, and their explicit image/retraction
obligations.  These are interfaces for deformation or flow proofs.
\end{description}

\section{Complex dynamics, Green functions, and puzzles}

\begin{description}
\item[\textbf{GreenSublevel}]
The module \lean{Mlc/Quadratic/Complex/GreenSublevel.lean} defines
\lean{Quadratic.GreenSublevel} and proves that it contains the
critical point for $c\in\M$.

\item[\lean{Mlc/GreenSublevelConnectedDirect.lean}]
Proves connectedness of the dynamical Green sublevels by direct potential
theory.

\item[\lean{Mlc/ParaPuzzleConnectivity.lean}]
Defines the set-level Green/intersection connectivity input, proves its
equivalence with the categorical version, and proves the conditional
parapuzzle connectedness theorem.

\item[\lean{Mlc/GreenSublevelIntersectionCounterexample.lean}]
Contains the certified $c=i$, $n=16$ obstruction.  The theorem
\lean{MLC.not_greenSublevelIntersectionCategoricalData} prevents the
full-sublevel model from being mistaken for the Yoccoz graph-cut model.

\item[\lean{Mlc/Quadratic/Complex/YoccozConformal.lean}]
Proves the conformal Yoccoz shrinking theorem from divergence of puzzle
annulus moduli:
\lean{Quadratic.yoccoz_theorem_conformal}.

\item[\lean{Mlc/Quadratic/Complex/ParaPuzzleBasis.lean}]
Shows how connected shrinking parameter pieces produce a neighbourhood basis
in the Mandelbrot subspace.
\end{description}

\section{Renormalization and current case interfaces}

\begin{description}
\item[\lean{Mlc/RenormalizationTypes.lean}]
Defines primitive, finite/infinite renormalization predicates and the
modulus-summability interface.

\item[\lean{Mlc/MoleculeRenormalizationTower.lean}]
Defines \lean{RenormalizationTower}, chooses its relations, and proves
monotonicity of cumulative periods.

\item[\lean{Mlc/SatelliteRenormalizationTower.lean}]
Specializes towers to satellite-renormalizable parameters.

\item[\lean{Mlc/MoleculeToParameterShrink.lean}]
Proves the principal-nest finite-union decomposition and the bridge from
uniform principal-nest bounds to shrinking parameter puzzles.

\item[\lean{Mlc/Literature/Combinatorics.lean}]
Adds explicit bounded-type, period-doubling, primitive/satellite/neutral
profile predicates and the arithmetic divergence theorem
\lean{not_summable_of_uniform_lower_bound}.

\item[\lean{Mlc/Literature/Renormalization.lean}]
Packages the a priori-bound certificate and proves its pointwise MLC
conclusion once connected parameter pieces are supplied.

\item[\lean{Mlc/Literature/Dudko.lean}]
Provides source-specific contracts for bounded-type satellite, Feigenbaum,
parabolically bounded primitive, and neutral-renormalization classes.
\end{description}

\section{Categorical and Efimov-facing interfaces}

\begin{description}
\item[\lean{Mlc/CategoricalTopologicalApproximation.lean}]
Implements the over-category of topological approximations, pullbacks, image
intersection, and factorization through subset inclusions.

\item[\lean{Mlc/EfimovCategoricalBridge.lean}]
Defines the native Mathlib-level substitutes for Pacman towers, strong
Mittag--Leffler data, graded invariant shadows, topological realizations,
and connected surjective morphisms.  It proves connectedness transport but
does not implement stable infinity-categories or continuous $K$-theory.

\item[\lean{Mlc/UniformGeometricRoot.lean}]
Assembles the conditional uniform retraction theorem and its categorical
form.

\item[\lean{Mlc/Literature/Categorical.lean}]
Provides review-facing wrappers:
\lean{categoricalMLC_of_rootInput},
\lean{categoricalMLC_of_finiteComponentCriterion},
\lean{categoricalMLC_of_squareRetractionTower}, and the explicit refutation
\lean{frozenGreenSublevelCarving_is_refuted}.
\end{description}

\section{Certified approximations, classification, and analytic support}

\begin{description}
\item[\lean{Mlc/CertifiedOrbitApproximation.lean}]
Defines finite-cell outer systems, compact stage unions, nested addresses,
and the diameter-control interface.  Its baseline and dyadic constructions
are sound outer approximations; they do not provide the missing local
connectivity buffer.

\item[\lean{Mlc/CertifiedTrappingRegions.lean}]
Defines finite trapping certificates for parameter boxes and proves their
soundness: a certified box is contained in $\M$.

\item[\lean{Mlc/FiniteTrappingRegions.lean}]
Extends the single-box certificate to finite unions and supplies the
period-two interior example used by the regression checks.

\item[\lean{Mlc/TrappingRegionObstruction.lean}]
Proves that the single-rectangle certificate scheme cannot be dense in
$\M$.  This is a negative result about one certificate design, not a
negative result about finite-union certificates.

\item[\lean{Mlc/ParameterAddressSpace.lean}]
Defines compatible dyadic addresses, nested intersections, and the
singleton/coverage conditions \lean{VD} and \lean{AC}.  The dyadic system
proves both conditions, but these address facts alone do not imply MLC.

\item[\lean{Mlc/ParameterClassification.lean}]
Proves the escape criterion and first-escape time, then gives an exhaustive
priority partition into attracting, parabolic, Siegel, Cremer, eventually
periodic, and residual predicates.  The residual class is deliberately
retained rather than declared empty.

\item[\lean{Mlc/ModelRegression.lean}]
Collects frozen-model exclusions and protects the project from silently
reintroducing a refuted Green-sublevel or trapping-region shortcut.
\end{description}

\section{Complex-analytic support modules}

The Böttcher and Green-function layer is split so that each analytic
construction has a small theorem surface:
\begin{itemize}
\item \lean{Mlc/Quadratic/Complex/Bottcher/LogSeries.lean} contains the
      logarithmic series and convergence estimates;
\item \lean{Mlc/Quadratic/Complex/Bottcher/CpowSlit.lean} handles the slit
      power map used in the coordinate construction;
\item \lean{Mlc/Quadratic/Complex/Bottcher/ConstructiveBasinCoordinate.lean}
      packages the basin coordinate;
\item \lean{Mlc/Quadratic/Complex/Bottcher/BottcherOnMDefs.lean} records
      the coordinate definitions and normalization;
\item \lean{Mlc/Quadratic/Complex/GreenHarmonic.lean} and
      \lean{Mlc/Quadratic/Complex/HarmonicMinimumPrinciple.lean} provide
      harmonicity and minimum-principle tools;
\item \lean{Mlc/Quadratic/Complex/GreenSublevel.lean} and
      \lean{Mlc/Quadratic/Complex/ConformalGroetzsch.lean} connect Green
      sublevels with modulus arguments.
\end{itemize}
The formal review treats these as analytic foundations, not as an automatic
proof of boundary landing or parameter-plane realization.

\section{Uniform geometry, residual categories, and regressions}

\begin{description}
\item[\lean{Mlc/UniformGeometricDomain.lean}]
Defines the fixed square $D=[-2,2]^2$ and the
\lean{OrbitRetractionTower} contract.

\item[\lean{Mlc/UniformGeometricApproximation.lean}]
Proves uniform Cauchy estimates and the continuous limit of an admissible
tower.

\item[\lean{Mlc/RetractionLocalConnectivity.lean}]
Turns the limit into a retraction and the retraction into local
connectedness.

\item[\lean{Mlc/UniformGeometricRoot.lean}]
Assembles these lemmas into the conditional square-tower theorem.

\item[\lean{Mlc/CategoricalResidual.lean}]
Packages the residual pseudo-Siegel and virtual-molecule inputs as a
categorical product/limit interface.  The inputs remain open.

\item[\lean{Mlc/GreenSublevelIntersectionCounterexample.lean}]
Contains the checked integer certificate and the theorem-level refutation
of the universal frozen Green-intersection proposition.
\end{description}

\section*{Formalization status}
\addcontentsline{toc}{section}{Formalization status}

The current codebase therefore proves the following hierarchy:
\[
\begin{gathered}
\text{finite orbit identities, compactness, and elementary topology}\\
\Downarrow\\
\text{conditional buffer, finite-cover, retraction, and carving theorems}\\
\Downarrow\\
\text{pointwise MLC from explicit puzzle/renormalization certificates}\\
\Downarrow\\
\text{global MLC only if one of the open existence packages is supplied}.
\end{gathered}
\]
It does not contain a project-level axiom asserting MLC, and the categorical
Green-sublevel intersection proposal is explicitly refuted.  This is the
correct formal boundary relative to the cited literature: known published
cases are represented by source-specific contracts, while the analytic
estimates and their parameter-plane realization remain the obligations that
the future formalization must discharge.
